\chapter{交错双线性型的特殊性质}

\section{交错双线性型的约化}

\begin{theorem}{PID上交错双线性型的约化}{}
	设$A$是PID,$E$是$n$维自由$A$-模,$\varphi\in\Bil_{A}\left(E\right)$交错,则存在$E$的基$\left(e_{i}\right)_{i=1}^{n}$,满足$2r\leq n$的$r\in\N$和$A\setminus\left\{0\right\}$的族$\left(\alpha_{i}\right)_{i=1}^{2r}$,使得
	\begin{enumerate}
		\item $\forall 1\leq i\leq r\left(\varphi\left(e_{2i-1},e_{2i}\right)=\alpha_{i}\right)$;
		\item $\forall 1\leq i\leq r-1\left(\alpha_{i}\mid\alpha_{i+1}\right)$;
		\item $\forall 1\leq i,j\leq n\left(i\leq j\implies\varphi\left(e_{i},e_{j}\right)=0\right)$;
		\item 理想序列$\left(A\alpha_{i}\right)_{i=1}^{r}$由条件(1)-(3)唯一确定;
		\item $E^{\circ}=\langle e_{2r+1},\ldots,e_{n}\rangle$.
	\end{enumerate}
\end{theorem}

本节接下来的内容中,默认$F$是域,$V$是$n$维$F$-向量空间.

\begin{corollary}{域上的辛基的存在性}{}
	设$\varphi\in\Bil_{F}\left(V\right)$,则存在$V$的基$\left(e_{i}\right)_{i=1}^{n}$和满足$2r\leq n$的$r\in\N$,对$A\times A$的每个族$\left(\left(a_{i},b_{i}\right)\right)_{i=1}^{n}$有
	\begin{align*}
		\varphi\left(\sum\limits_{i=1}^{n}a_{i}e_{i},\sum\limits_{i=1}^{n}b_{i}e_{i}\right)=\sum\limits_{i=1}^{r}\left(a_{2i-1}b_{2i}-a_{2i}b_{2i-1}\right).
	\end{align*}
	特别地,$\rank\left(\varphi\right)=2r$.
\end{corollary}

\begin{definition}{双线性型的辛基}{}
	设$\left(e_{i}\right)_{i=1}^{n}$是$V$的基,$r\in\N$且$2r\leq n$.若对$A\times A$的每个族$\left(\left(a_{i},b_{i}\right)\right)_{i=1}^{n}$有
	\begin{align*}
		\varphi\left(\sum\limits_{i=1}^{n}a_{i}e_{i},\sum\limits_{i=1}^{n}b_{i}e_{i}\right)=\sum\limits_{i=1}^{r}\left(a_{2i-1}b_{2i}-a_{2i}b_{2i-1}\right),
	\end{align*}
	则称$\left(e_{i}\right)_{i=1}^{n}$是$V$的辛基.
\end{definition}

\begin{corollary}{$2$-向量的标准分解}{}
	对每个$z\in\bigwedge\limits^{2}V$,存在$V$的基$\left(e_{i}\right)_{i=1}^{n}$和$r\in\N$,使得$2r\leq n$且$z=\sum\limits_{i=1}^{r}\left(e_{2i-1}\wedge e_{2i}\right)$.
\end{corollary}

\begin{corollary}{交错矩阵的相抵标准型}{}
	设$\boldsymbol{R}\in\mathbf{M}_{n}\left(F\right)$交错,则存在$r\in\N$使得$\rank\left(\boldsymbol{A}\right)=2r$,并且存在$\boldsymbol{P}\in\GL_{n}\left(F\right)$使得
	\begin{align*}
		\boldsymbol{P}^{\top}\boldsymbol{A}\boldsymbol{P}=\begin{bmatrix}
			\boldsymbol{O}_{r\times r}&\boldsymbol{I}_{r\times r}&\boldsymbol{O}_{\left(n-2r\right)\times\left(n-2r\right)}\\
			-\boldsymbol{I}_{r\times r}&\boldsymbol{O}_{r\times r}&\boldsymbol{O}_{\left(n-2r\right)\times\left(n-2r\right)}\\
			\boldsymbol{O}_{\left(n-2r\right)\times\left(n-2r\right)}&\boldsymbol{O}_{\left(n-2r\right)\times r}&\boldsymbol{O}_{\left(n-2r\right)\times\left(n-2r\right)}
		\end{bmatrix}.
	\end{align*}
\end{corollary}

\begin{remark}{}{}
	如果$A$是交换环,那么对每个$n\in\N$和$A$上的每个$2n+1$阶方阵$\boldsymbol{A}$都有$\det\left(\boldsymbol{A}\right)=0$.
\end{remark}
